\documentclass[12pt,a4paper]{article}
\usepackage{ctex}
\usepackage{amsmath,amssymb}
\usepackage{graphicx}
\usepackage{listings}
\usepackage{xcolor}
\usepackage{geometry}
\usepackage{float}
\usepackage{caption}
\geometry{left=2.5cm,right=2.5cm,top=2.5cm,bottom=2.5cm}

% Verilog 代码高亮
\lstset{
    language=Verilog,
    basicstyle=\ttfamily\small,
    keywordstyle=\color{blue}\bfseries,
    commentstyle=\color{purple!60},
    stringstyle=\color{red},
    numbers=left,
    numberstyle=\tiny\color{gray},
    numbersep=5pt,
    frame=single,
    framerule=0.3pt,
    backgroundcolor=\color{white},
    breaklines=true,
    showstringspaces=false,
    tabsize=4
}

% XDC 约束高亮
\lstdefinelanguage{XDC}{
    keywords={set_property, PACKAGE_PIN, IOSTANDARD, get_ports},
    keywordstyle=\color{blue}\bfseries,
    commentstyle=\color{purple!60},
    stringstyle=\color{red},
    alsoletter={[,]},
}

\begin{document}
\title{\textbf{实验三 寄存器堆功能改进实验报告}}
\date{\today}
\maketitle

\section{实验目的}
\begin{enumerate}
    \item 理解 32 位通用寄存器堆的硬件结构与读写工作原理。
    \item 将写使能信号从 1 位扩展为 2 位，实现高低 16 位独立写入控制。
    \item 新增专用寄存器，实时计算并存储 1 号与 2 号寄存器数据之和。
    \item 完成 FPGA 板级验证，验证改进后寄存器堆功能正确性。
\end{enumerate}

\section{实验原理与改进方案}
\subsection{原有寄存器堆结构}
传统寄存器堆包含 32 个 32 位寄存器，两路异步读、一路同步写，使用 1 位写使能控制整个 32 位数据写入。

\subsection{本次改进内容}
\begin{enumerate}
    \item \textbf{写使能扩展为 2 位}
    \begin{itemize}
        \item \texttt{wen[0]}：控制写数据低 16 位 \texttt{wdata[15:0]}
        \item \texttt{wen[1]}：控制写数据高 16 位 \texttt{wdata[31:16]}
    \end{itemize}
    可实现只写高 16 位、只写低 16 位或完整 32 位写入。

    \item \textbf{新增专用求和寄存器}
    增加寄存器 \texttt{sum\_reg}，在时钟上升沿自动计算
    \[
    \text{sum\_reg} = \text{rf[1]} + \text{rf[2]}
    \]

    \item \textbf{调试端口扩展}
    将测试地址 \texttt{test\_addr = 31} 映射为 \texttt{sum\_reg}，可直接观察求和结果。
\end{enumerate}

\subsection{改进后寄存器堆原理图说明}
\begin{figure}[H]
    \centering
    \includegraphics[width=0.8\textwidth]{3.jpg}
    \caption{改进后寄存器堆硬件原理图}
\end{figure}

原理图结构说明：
\begin{enumerate}
    \item 写控制逻辑：2 位写使能分别控制高低 16 位写入使能。
    \item 加法模块：时钟同步加法器，实时计算 rf[1]+rf[2]。
    \item 所有写操作与求和操作均在时钟上升沿同步完成。
\end{enumerate}

\section{改进后 Verilog 代码实现}
\subsection{寄存器堆核心模块 regfile.v}
\begin{lstlisting}
`timescale 1ns / 1ps
module regfile(
    input             clk,
    input      [1:0]  wen, // 2位写使能
    input      [4 :0] raddr1,
    input      [4 :0] raddr2,
    input      [4 :0] waddr,
    input      [31:0] wdata,
    output reg [31:0] rdata1,
    output reg [31:0] rdata2,
    input      [4 :0] test_addr,
    output reg [31:0] test_data
);
    reg [31:0] rf[31:0];
    reg [31:0] sum_reg;  // 新增求和寄存器

    // 分高低16位写入
    always @(posedge clk)
    begin
        if(waddr != 5'd0) begin
            if(wen[0]) rf[waddr][15:0]  <= wdata[15:0];
            if(wen[1]) rf[waddr][31:16] <= wdata[31:16];
        end
        sum_reg <= rf[1] + rf[2];
    end

    always @(*)
    begin
        case (raddr1)
            5'd1 : rdata1 <= rf[1 ];
            5'd2 : rdata1 <= rf[2 ];
            5'd3 : rdata1 <= rf[3 ];
            5'd4 : rdata1 <= rf[4 ];
            5'd5 : rdata1 <= rf[5 ];
            5'd6 : rdata1 <= rf[6 ];
            5'd7 : rdata1 <= rf[7 ];
            5'd8 : rdata1 <= rf[8 ];
            5'd9 : rdata1 <= rf[9 ];
            5'd10: rdata1 <= rf[10];
            5'd11: rdata1 <= rf[11];
            5'd12: rdata1 <= rf[12];
            5'd13: rdata1 <= rf[13];
            5'd14: rdata1 <= rf[14];
            5'd15: rdata1 <= rf[15];
            5'd16: rdata1 <= rf[16];
            5'd17: rdata1 <= rf[17];
            5'd18: rdata1 <= rf[18];
            5'd19: rdata1 <= rf[19];
            5'd20: rdata1 <= rf[20];
            5'd21: rdata1 <= rf[21];
            5'd22: rdata1 <= rf[22];
            5'd23: rdata1 <= rf[23];
            5'd24: rdata1 <= rf[24];
            5'd25: rdata1 <= rf[25];
            5'd26: rdata1 <= rf[26];
            5'd27: rdata1 <= rf[27];
            5'd28: rdata1 <= rf[28];
            5'd29: rdata1 <= rf[29];
            5'd30: rdata1 <= rf[30];
            5'd31: rdata1 <= rf[31];
            default : rdata1 <= 32'd0;
        endcase
    end

    always @(*)
    begin
        case (raddr2)
            5'd1 : rdata2 <= rf[1 ];
            5'd2 : rdata2 <= rf[2 ];
            5'd3 : rdata2 <= rf[3 ];
            5'd4 : rdata2 <= rf[4 ];
            5'd5 : rdata2 <= rf[5 ];
            5'd6 : rdata2 <= rf[6 ];
            5'd7 : rdata2 <= rf[7 ];
            5'd8 : rdata2 <= rf[8 ];
            5'd9 : rdata2 <= rf[9 ];
            5'd10: rdata2 <= rf[10];
            5'd11: rdata2 <= rf[11];
            5'd12: rdata2 <= rf[12];
            5'd13: rdata2 <= rf[13];
            5'd14: rdata2 <= rf[14];
            5'd15: rdata2 <= rf[15];
            5'd16: rdata2 <= rf[16];
            5'd17: rdata2 <= rf[17];
            5'd18: rdata2 <= rf[18];
            5'd19: rdata2 <= rf[19];
            5'd20: rdata2 <= rf[20];
            5'd21: rdata2 <= rf[21];
            5'd22: rdata2 <= rf[22];
            5'd23: rdata2 <= rf[23];
            5'd24: rdata2 <= rf[24];
            5'd25: rdata2 <= rf[25];
            5'd26: rdata2 <= rf[26];
            5'd27: rdata2 <= rf[27];
            5'd28: rdata2 <= rf[28];
            5'd29: rdata2 <= rf[29];
            5'd30: rdata2 <= rf[30];
            5'd31: rdata2 <= rf[31];
            default : rdata2 <= 32'd0;
        endcase
    end

    always @(*)
    begin
        case (test_addr)
            5'd1 : test_data <= rf[1 ];
            5'd2 : test_data <= rf[2 ];
            5'd3 : test_data <= rf[3 ];
            5'd4 : test_data <= rf[4 ];
            5'd5 : test_data <= rf[5 ];
            5'd6 : test_data <= rf[6 ];
            5'd7 : test_data <= rf[7 ];
            5'd8 : test_data <= rf[8 ];
            5'd9 : test_data <= rf[9 ];
            5'd10: test_data <= rf[10];
            5'd11: test_data <= rf[11];
            5'd12: test_data <= rf[12];
            5'd13: test_data <= rf[13];
            5'd14: test_data <= rf[14];
            5'd15: test_data <= rf[15];
            5'd16: test_data <= rf[16];
            5'd17: test_data <= rf[17];
            5'd18: test_data <= rf[18];
            5'd19: test_data <= rf[19];
            5'd20: test_data <= rf[20];
            5'd21: test_data <= rf[21];
            5'd22: test_data <= rf[22];
            5'd23: test_data <= rf[23];
            5'd24: test_data <= rf[24];
            5'd25: test_data <= rf[25];
            5'd26: test_data <= rf[26];
            5'd27: test_data <= rf[27];
            5'd28: test_data <= rf[28];
            5'd29: test_data <= rf[29];
            5'd30: test_data <= rf[30];
            5'd31: test_data <= sum_reg; // 设置专用寄存器
            default : test_data <= 32'd0;
        endcase
    end
endmodule
\end{lstlisting}

\subsection{显示驱动模块 regfile\_display.v}
\begin{lstlisting}
module regfile_display(
    input clk,
    input resetn,

    input [1:0] wen, // 2位写使能
    input [1:0] input_sel,

    output led_wen,
    output led_waddr,
    output led_wdata,
    output led_raddr1,
    output led_raddr2,

    output lcd_rst,
    output lcd_cs,
    output lcd_rs,
    output lcd_wr,
    output lcd_rd,
    inout[15:0] lcd_data_io,
    output lcd_bl_ctr,
    inout ct_int,
    inout ct_sda,
    output ct_scl,
    output ct_rstn
);
    assign led_wen    = |wen; // 控制写使能 LED 指示灯正常工作
    assign led_raddr1 = (input_sel==2'd0);
    assign led_raddr2 = (input_sel==2'd1);
    assign led_waddr  = (input_sel==2'd2);
    assign led_wdata  = (input_sel==2'd3);

    wire [31:0] test_data;
    wire [4 :0] test_addr;
    reg  [4 :0] raddr1;
    reg  [4 :0] raddr2;
    reg  [4 :0] waddr;
    reg  [31:0] wdata;
    wire [31:0] rdata1;
    wire [31:0] rdata2;

    regfile rf_module(
        .clk      (clk     ),
        .wen      (wen     ),
        .raddr1   (raddr1  ),
        .raddr2   (raddr2  ),
        .waddr    (waddr   ),
        .wdata    (wdata   ),
        .rdata1   (rdata1  ),
        .rdata2   (rdata2  ),
        .test_addr(test_addr),
        .test_data(test_data)
    );

    reg         display_valid;
    reg  [39:0] display_name;
    reg  [31:0] display_value;
    wire [5 :0] display_number;
    wire        input_valid;
    wire [31:0] input_value;

    lcd_module lcd_module(
        .clk            (clk           ),
        .resetn         (resetn        ),
        .display_valid  (display_valid ),
        .display_name   (display_name  ),
        .display_value  (display_value ),
        .display_number (display_number),
        .input_valid    (input_valid   ),
        .input_value    (input_value   ),
        .lcd_rst        (lcd_rst       ),
        .lcd_cs         (lcd_cs        ),
        .lcd_rs         (lcd_rs        ),
        .lcd_wr         (lcd_wr        ),
        .lcd_rd         (lcd_rd        ),
        .lcd_data_io    (lcd_data_io   ),
        .lcd_bl_ctr     (lcd_bl_ctr    ),
        .ct_int         (ct_int        ),
        .ct_sda         (ct_sda        ),
        .ct_scl         (ct_scl        ),
        .ct_rstn        (ct_rstn       )
    );

    assign test_addr = display_number - 5'd7;

    always @(posedge clk) begin
        if (!resetn) raddr1 <= 5'd0;
        else if (input_valid && input_sel == 2'd0) raddr1 <= input_value[4:0];
    end

    always @(posedge clk) begin
        if (!resetn) raddr2 <= 5'd0;
        else if (input_valid && input_sel == 2'd1) raddr2 <= input_value[4:0];
    end

    always @(posedge clk) begin
        if (!resetn) waddr <= 5'd0;
        else if (input_valid && input_sel == 2'd2) waddr <= input_value[4:0];
    end

    always @(posedge clk) begin
        if (!resetn) wdata <= 32'd0;
        else if (input_valid && input_sel == 2'd3) wdata <= input_value;
    end

    always @(posedge clk) begin
        if (display_number > 6'd6 && display_number < 6'd39) begin
            display_valid <= 1'b1;
            display_name[39:16] <= "REG";
            display_name[15: 8] <= {4'b0011,3'b000,test_addr[4]};
            display_name[7 : 0] <= {4'b0011,test_addr[3:0]};
            display_value <= test_data;
        end
        else begin
            case(display_number)
                6'd1: begin display_valid<=1; display_name<="RADD1"; display_value<=raddr1; end
                6'd2: begin display_valid<=1; display_name<="RDAT1"; display_value<=rdata1; end
                6'd3: begin display_valid<=1; display_name<="RADD2"; display_value<=raddr2; end
                6'd4: begin display_valid<=1; display_name<="RDAT2"; display_value<=rdata2; end
                6'd5: begin display_valid<=1; display_name<="WADDR"; display_value<=waddr;  end
                6'd6: begin display_valid<=1; display_name<="WDATA"; display_value<=wdata;  end
                default: begin display_valid<=0; display_name<=0; display_value<=0; end
            endcase
        end
    end
endmodule
\end{lstlisting}

\subsection{管脚约束文件}
\begin{lstlisting}[language=XDC]
#时钟信号连接
set_property PACKAGE_PIN AC19 [get_ports clk]

#复位按键，低电平有效
set_property PACKAGE_PIN Y3 [get_ports resetn]

#LED指示灯
set_property PACKAGE_PIN H7 [get_ports led_wen]
set_property PACKAGE_PIN D5 [get_ports led_waddr]
set_property PACKAGE_PIN A3 [get_ports led_wdata]
set_property PACKAGE_PIN A5 [get_ports led_raddr1]
set_property PACKAGE_PIN A4 [get_ports led_raddr2]

#拨码开关 (wen[0] 和 wen[1] 分别分配到具体的硬件引脚)
set_property PACKAGE_PIN AC21 [get_ports wen[0]]
set_property PACKAGE_PIN AD24 [get_ports wen[1]] 
set_property PACKAGE_PIN AC22 [get_ports input_sel[0]]
set_property PACKAGE_PIN AC23 [get_ports input_sel[1]]

set_property IOSTANDARD LVCMOS33 [get_ports clk]
set_property IOSTANDARD LVCMOS33 [get_ports resetn]
set_property IOSTANDARD LVCMOS33 [get_ports led_wen]
set_property IOSTANDARD LVCMOS33 [get_ports led_raddr1]
set_property IOSTANDARD LVCMOS33 [get_ports led_raddr2]
set_property IOSTANDARD LVCMOS33 [get_ports led_waddr]
set_property IOSTANDARD LVCMOS33 [get_ports led_wdata]
set_property IOSTANDARD LVCMOS33 [get_ports wen]
set_property IOSTANDARD LVCMOS33 [get_ports input_sel[1]]
set_property IOSTANDARD LVCMOS33 [get_ports input_sel[0]]

#LCD接口
set_property PACKAGE_PIN J25 [get_ports lcd_rst]
set_property PACKAGE_PIN H18 [get_ports lcd_cs]
set_property PACKAGE_PIN K16 [get_ports lcd_rs]
set_property PACKAGE_PIN L8 [get_ports lcd_wr]
set_property PACKAGE_PIN K8 [get_ports lcd_rd]
set_property PACKAGE_PIN J15 [get_ports lcd_bl_ctr]
set_property PACKAGE_PIN H9 [get_ports {lcd_data_io[0]}]
set_property PACKAGE_PIN K17 [get_ports {lcd_data_io[1]}]
set_property PACKAGE_PIN J20 [get_ports {lcd_data_io[2]}]
set_property PACKAGE_PIN M17 [get_ports {lcd_data_io[3]}]
set_property PACKAGE_PIN L17 [get_ports {lcd_data_io[4]}]
set_property PACKAGE_PIN L18 [get_ports {lcd_data_io[5]}]
set_property PACKAGE_PIN L15 [get_ports {lcd_data_io[6]}]
set_property PACKAGE_PIN M15 [get_ports {lcd_data_io[7]}]
set_property PACKAGE_PIN M16 [get_ports {lcd_data_io[8]}]
set_property PACKAGE_PIN L14 [get_ports {lcd_data_io[9]}]
set_property PACKAGE_PIN M14 [get_ports {lcd_data_io[10]}]
set_property PACKAGE_PIN F22 [get_ports {lcd_data_io[11]}]
set_property PACKAGE_PIN G22 [get_ports {lcd_data_io[12]}]
set_property PACKAGE_PIN G21 [get_ports {lcd_data_io[13]}]
set_property PACKAGE_PIN H24 [get_ports {lcd_data_io[14]}]
set_property PACKAGE_PIN J16 [get_ports {lcd_data_io[15]}]
set_property PACKAGE_PIN L19 [get_ports ct_int]
set_property PACKAGE_PIN J24 [get_ports ct_sda]
set_property PACKAGE_PIN H21 [get_ports ct_scl]
set_property PACKAGE_PIN G24 [get_ports ct_rstn]

set_property IOSTANDARD LVCMOS33 [get_ports lcd_rst]
set_property IOSTANDARD LVCMOS33 [get_ports lcd_cs]
set_property IOSTANDARD LVCMOS33 [get_ports lcd_rs]
set_property IOSTANDARD LVCMOS33 [get_ports lcd_wr]
set_property IOSTANDARD LVCMOS33 [get_ports lcd_rd]
set_property IOSTANDARD LVCMOS33 [get_ports lcd_bl_ctr]
set_property IOSTANDARD LVCMOS33 [get_ports {lcd_data_io[0]}]
set_property IOSTANDARD LVCMOS33 [get_ports {lcd_data_io[1]}]
set_property IOSTANDARD LVCMOS33 [get_ports {lcd_data_io[2]}]
set_property IOSTANDARD LVCMOS33 [get_ports {lcd_data_io[3]}]
set_property IOSTANDARD LVCMOS33 [get_ports {lcd_data_io[4]}]
set_property IOSTANDARD LVCMOS33 [get_ports {lcd_data_io[5]}]
set_property IOSTANDARD LVCMOS33 [get_ports {lcd_data_io[6]}]
set_property IOSTANDARD LVCMOS33 [get_ports {lcd_data_io[7]}]
set_property IOSTANDARD LVCMOS33 [get_ports {lcd_data_io[8]}]
set_property IOSTANDARD LVCMOS33 [get_ports {lcd_data_io[9]}]
set_property IOSTANDARD LVCMOS33 [get_ports {lcd_data_io[10]}]
set_property IOSTANDARD LVCMOS33 [get_ports {lcd_data_io[11]}]
set_property IOSTANDARD LVCMOS33 [get_ports {lcd_data_io[12]}]
set_property IOSTANDARD LVCMOS33 [get_ports {lcd_data_io[13]}]
set_property IOSTANDARD LVCMOS33 [get_ports {lcd_data_io[14]}]
set_property IOSTANDARD LVCMOS33 [get_ports {lcd_data_io[15]}]
set_property IOSTANDARD LVCMOS33 [get_ports ct_int]
set_property IOSTANDARD LVCMOS33 [get_ports ct_sda]
set_property IOSTANDARD LVCMOS33 [get_ports ct_scl]
set_property IOSTANDARD LVCMOS33 [get_ports ct_rstn]
\end{lstlisting}

\section{实验步骤}
\begin{enumerate}
    \item 新建 Vivado 工程，添加上述 Verilog 设计文件与约束文件。
    \item 完成综合、实现、比特流生成。
    \item 将比特流下载至 FPGA 实验箱。
    \item 通过拨码开关与按键配置读写地址、数据与写使能。
    \item 观察 LCD 显示结果，验证高低 16 位独立写入与求和功能。
\end{enumerate}

\section{实验箱验证结果与说明}
\begin{figure}[H]
    \centering
    \includegraphics[width=0.8\textwidth]{1.jpg}
    \caption{添加求和操作}
\end{figure}

\begin{figure}[H]
    \centering
    \includegraphics[width=0.8\textwidth]{2.jpg}
    \caption{改进写操作}
\end{figure}

\subsection{验证内容与数据说明}
\begin{enumerate}
    \item \textbf{专用求和寄存器测试}
    \begin{itemize}
        \item \texttt{REG01 + REG02 = REG1F (1 + 1 = 2)}
        \item 与理论计算一致，求和功能正确。
    \end{itemize}

    \item \textbf{高低 16 位独立写入测试(在实验1基础上进行测试)}
    \begin{itemize}
        \item 写数据：\texttt{wdata = 0x7FFFFFFF}
        \item 仅使能 \texttt{wen[0]=1}，写入低 16 位，REG02 显示为 0x0000FFFF
        \item 仅使能 \texttt{wen[1]=1}，写入高 16 位，REG01 显示为 0x7FFF0001
        \item 经实验可得，独立写入功能正确
    \end{itemize}

\end{enumerate}

\section{实验总结}
\begin{enumerate}
    \item 成功将寄存器堆写使能扩展为 2 位，实现高低 16 位独立可控写入。
    \item 实现了专用求和寄存器，可实时计算并保存 1 号与 2 号寄存器之和。
    \item 完成 FPGA 板级验证，所有功能均正常工作，符合设计要求。
    \item 进一步理解了寄存器堆的内部结构、同步时序与多路控制逻辑。
\end{enumerate}

\end{document}